\section{Decisiones tecnológicas}

\begin{tabularx}{\linewidth}{@{}p{0.15\linewidth} p{0.24\linewidth} Y@{}}
\toprule
\textbf{Componente} & \textbf{Elección final} & \textbf{Justificación} \\
\midrule
Hipervisor & Proxmox VE 9.2.2; kernel 7.0.2-6-pve & Hipervisor Tipo 1 con administración web y soporte KVM/QEMU. \\
Cliente & openSUSE Leap 16.0; KDE Plasma 6.4.2 & Distribución distinta de Ubuntu con navegador, terminal y FileZilla. \\
Servidores & Ubuntu Server 26.04.1 LTS & Versión exigida; \texttt{multi-user.target} evita entorno gráfico. \\
Web & Apache 2.4.66, PHP 8.5.4 y WordPress 7.1.4 & Versiones estables de Ubuntu 26.04 y WordPress disponibles al instalar; VirtualHosts separan HTTP y HTTPS \cite{ubuntu-apache,ubuntu-php,wordpress-releases}. \\
Base de datos & MySQL 8.4.11, puerto 13306 & Versión LTS distribuida para Ubuntu 26.04; separa persistencia y evita el puerto predeterminado \cite{ubuntu-mysql}. \\
DNS & Amazon Route 53 & Zona pública autoritativa para \texttt{laboratorio5.com}. \\
Certificado & Let's Encrypt, DNS-01 & Valida el dominio sin publicar la IP privada del servidor \cite{letsencrypt-challenges}. \\
Firewall & UFW & Reglas simples, auditables y restringidas por IP origen. \\
\bottomrule
\end{tabularx}

\subsection{Registro y comprobación de versiones}

Las versiones se registraron desde los sistemas instalados y se contrastaron con las fuentes oficiales de los paquetes y de WordPress. No se dedujeron a partir del nombre de una ISO. Los comandos de comprobación se muestran para que el registro sea reproducible.

\begin{lstlisting}[caption={Comandos utilizados para comprobar versiones.}]
# Host Proxmox
pveversion
uname -r

# Cliente openSUSE
cat /etc/os-release
plasmashell --version

# Servidor web
cat /etc/os-release
apache2 -v
php -v
php -r 'include "/var/www/wordpress/wp-includes/version.php";
echo $wp_version, PHP_EOL;'

# Servidor DB
cat /etc/os-release
mysql --version
sudo mysql -NBe 'SELECT VERSION();'
\end{lstlisting}

El resultado consolidado fue: Proxmox VE 9.2.2 con kernel 7.0.2-6-pve; openSUSE Leap 16.0 con KDE Plasma 6.4.2; Ubuntu Server 26.04.1 LTS en ambos servidores; Apache 2.4.66; PHP 8.5.4; WordPress 7.1.4; y MySQL 8.4.11. Las revisiones de paquete de Ubuntu pueden cambiar con actualizaciones de seguridad, por lo que el informe conserva la versión funcional relevante y no depende del sufijo de empaquetado.

\subsection{Principios aplicados}

Se aplicaron separación de funciones, mínimo privilegio y defensa en profundidad. La base de datos no está disponible para el cliente; el acceso SSH requiere llaves; \texttt{root} no ingresa remotamente; los archivos de WordPress solo son modificables por \texttt{wpadmin}; el firewall refuerza las restricciones de escucha de cada servicio.

WordPress recomienda HTTPS, PHP y MySQL en versiones mantenidas \cite{wordpress-requirements}. La solución evita publicar una dirección RFC 1918 en DNS público: Route 53 se usa únicamente para el desafío TXT y \texttt{/etc/hosts} resuelve el dominio dentro del laboratorio.
